\section{本章小结}
\label{sec:lde-summary}

线性方程把初始状态和外部输入线性地映成响应，本章的所有公式都是这个线性映射的不同写法。标量方程用积分因子得到显式解，公式在系数连续的整个区间上成立；系数只可积时，同一公式在积分形式下仍然成立，解是绝对连续函数。

常系数系统的传播是矩阵指数，它满足齐次方程和单位初值。对角化把系统分解成独立的模式，复特征值对应旋转和伸缩，三角形式让系统逐行回代。特征根决定解的成分，但只看特征值是不够的：相同的特征值可以有不同数量的特征向量，广义特征向量带来$t$的多项式因子和有限时间的放大。高阶标量方程经状态降阶化为友矩阵系统，重根必定带来多项式因子。

时变系统的传播是状态转移矩阵。它依赖初始时刻和终止时刻，满足复合关系，由基本解矩阵算出。各时刻的系数不可交换时，时间次序不能丢弃，$\exp\int\bm{A}$不是传播；所有时刻的系数都可交换时，这个简化才成立。

非齐次方程的解等于零输入响应加零状态响应，常数变易公式把输入看成不断注入、并继续传播的状态增量。共振来自输入与齐次模式的配合，它与Jordan链上的多项式因子是同一机制。因果Green函数是单位脉冲的响应，它在输入发生之前为零，离散以后是下三角矩阵。Laplace变换把常系数方程化为代数方程，初值自动进入，极点的位置和重数对应解的增长、振荡和多项式因子。

表~\ref{tab:lde-methods}对照了这些方法的适用范围。

\begin{BookInlineTable}
  \centering
  \caption{线性初值问题的四种工具。}
  \label{tab:lde-methods}
  \begin{tabular}{lll}
    \toprule
    方法 & 系数 & 输入\\
    \midrule
    积分因子 & 标量，连续或局部可积 & 连续或局部可积\\
    矩阵指数 & 常矩阵 & 连续或局部可积\\
    状态转移矩阵 & 连续或局部可积的矩阵 & 连续或局部可积\\
    Laplace变换 & 常系数 & 指数阶、分段连续\\
    \bottomrule
  \end{tabular}
\end{BookInlineTable}

非线性方程在一条已知轨迹附近的小扰动，满足的正是线性时变方程，所以状态转移矩阵和常数变易公式会在第~\ref{chap:ode-sensitivity}章重新出现。本章的结论都是有限时间上的显式刻画；响应在无穷长时间内是否稳定、输入按频率怎样被放大，以及怎样设计反馈来改变传播，都不是这里的问题。
